% web: the page: <head>, MathJax, fonts, colors and layout.
\makeatletter

% --- <head> -----------------------------------------------------------------

% values from TeX for the Lua passes, as <meta name="tex-.."> (web.mk4 reads
% and removes them): the size of the math extension font as a percentage of
% the body font (see delimiters.extension_scale), and the baselineskip in
% ems (see tables.lua)
\setbox\z@\hbox{\everymath{}$\xdef\web@exscale{\web@ratio{100\fontdimen6\textfont3}}$}
\xdef\web@baselineskip{\web@ratio\baselineskip}
% MathJax rejects -tex-calligraphic (\mathcal) unless told not to check, and
% its config has to load before MathJax, so rebuild <head> with the config
% first. The config also picks MathJax's Computer Modern font (its default,
% New Computer Modern, is heavier than the PDF, and its Latin Modern
% misplaces accents on Greek letters), scrolls display math that doesn't
% fit, and sets math at the size of the text, as TeX does (rather than
% matching x-heights, which comes out a bit small)
\Configure{@HEAD}{}
\Configure{@HEAD}{\HCode{<meta charset="utf-8" />\Hnewline}}
\Configure{@HEAD}{\HCode{<meta name="viewport" content="width=device-width,initial-scale=1" />\Hnewline}}
\Configure{@HEAD}{\HCode{<link rel="stylesheet" type="text/css" href="\jobname.css" />\Hnewline}}
\Configure{@HEAD}{\HCode{<meta name="tex-exscale" content="\web@exscale" />\Hnewline}}
\Configure{@HEAD}{\HCode{<meta name="tex-baselineskip" content="\web@baselineskip" />\Hnewline}}
\Configure{@HEAD}{\HCode{<script>window.MathJax = {
    mml: { verify: { checkMathvariants: false } },
    output: { font: 'mathjax-tex', displayOverflow: 'scroll',
        matchFontHeight: false } };</script>\Hnewline}}
\Configure{@HEAD}{\HCode{<script async="async" id="MathJax-script" src="https://cdn.jsdelivr.net/npm/mathjax@4/mml-chtml.js"></script>\Hnewline}}

% --- colors -----------------------------------------------------------------

% The page follows the reader's light or dark setting, unless the root
% element has data-theme="light" or "dark". Light mode uses the PDF's colors.
% \WebColor{name}{light color}{dark color} (xcolor colors) makes the CSS
% variable --name, which is one or the other
\newcommand\WebColor[3]{%
    \ifdefined\convertcolorspec
        \web@hex{#2}\web@light
        \web@hex{#3}\web@dark
        \web@Css{:root { --#1: light-dark(\#\web@light, \#\web@dark); }}%
    \else
        \PackageWarning{web}{\string\WebColor\space needs xcolor}%
    \fi
}
\newcommand\web@hex[2]{% color, macro for its hex code
    \extractcolorspecs{#1}\web@model\web@spec
    \convertcolorspec\web@model\web@spec{HTML}#2%
}
\Css{:root { color-scheme: light dark; }}
\Css{:root[data-theme="light"] { color-scheme: light; }}
\Css{:root[data-theme="dark"] { color-scheme: dark; }}
\Css{:root { --web-bg: light-dark(\#FFFFFF, \#1E1E1E); --web-text: light-dark(\#000000, \#E3E3E3); }}
% (html body: to win over tex4ht's own body rules)
\Css{html body { background-color: var(--web-bg); color: var(--web-text); }}
% links: hyperref's colors with colorlinks (otherwise the PDF draws boxes,
% and the links here keep the color of the text). In a macro, so TeX doesn't
% meet \ifHy@colorlinks when hyperref isn't loaded
\newcommand\web@linkcolors{%
    \ifHy@colorlinks
        \definecolor{web-link-dark}{RGB}{108, 180, 255}%
        \definecolor{web-cite-dark}{RGB}{95, 211, 160}%
        \definecolor{web-url-dark}{RGB}{242, 127, 176}%
        \WebColor{web-link}{\@linkcolor}{web-link-dark}%
        \WebColor{web-cite}{\@citecolor}{web-cite-dark}%
        \WebColor{web-url}{\@urlcolor}{web-url-dark}%
    \fi
}
\@ifpackageloaded{hyperref}{\web@linkcolors}{}
\Css{a { color: var(--web-link, inherit); text-decoration: none; }}
% (tex4ht names bibliography anchors X..)
\Css{a[href\string^="\#X"] { color: var(--web-cite, inherit); }}
\Css{a.url, a[href\string^="mailto:"] { color: var(--web-url, inherit); }}
% pictures tex4ht generates (TikZ; \jobname0x.svg and so on) are black line
% art; invert them in dark mode
\Css{@media (prefers-color-scheme: dark) {
    :root:not([data-theme="light"]) { --web-picture-filter: invert(1); } }}
\Css{:root[data-theme="dark"] { --web-picture-filter: invert(1); }}
\Css{img[src\string^="\jobname"][src\string$="x.svg"] { filter: var(--web-picture-filter, none); }}

% --- fonts ------------------------------------------------------------------

% Latin Modern, which web.mk4 copies into fonts/
\Css{@font-face { font-family: "Latin Modern Roman";
    src: url("fonts/lmroman10-regular.otf"); }}
\Css{@font-face { font-family: "Latin Modern Roman"; font-weight: bold;
    src: url("fonts/lmroman10-bold.otf"); }}
\Css{@font-face { font-family: "Latin Modern Roman"; font-style: italic;
    src: url("fonts/lmroman10-italic.otf"); }}
\Css{@font-face { font-family: "Latin Modern Roman"; font-weight: bold; font-style: italic;
    src: url("fonts/lmroman10-bolditalic.otf"); }}
\Css{@font-face { font-family: "Latin Modern Mono";
    src: url("fonts/lmmono10-regular.otf"); }}
\Css{html body { font-family: "Latin Modern Roman", serif; }}
% typewriter text, in tex4ht's class names for Latin Modern and for
% Computer Modern (without lmodern)
\Css{[class*="-lmtt"], [class\string^="cmtt-"], [class\string^="ectt-"] { font-family: "Latin Modern Mono", monospace; }}
% Latin Modern's no-break space (~) is 0.5em wide, its space 0.333em;
% text.no_break_spaces marks no-break spaces, and keeps one together with
% the math after it
\Css{.nbsp { letter-spacing: -0.167em; }}
\Css{.nowrap { white-space: nowrap; }}

% --- layout -----------------------------------------------------------------

% Text at 12pt, larger than most PDFs, for reading on a screen, and
% everything else in proportion: headings as much larger than the text as
% in the PDF (in the standard classes' sizes); MathJax sets math as large as
% the text around it, and tex4ht sizes other fonts relative to the text
\Css{html body { font-size: 12pt; }}
\newcommand\web@headsize[2]{% selector, size command
    \begingroup#2\xdef\web@size{\f@size pt}\endgroup
    \web@Css{#1 { font-size: \web@em\web@size; }}%
}
\web@headsize{h1.titleHead}{\LARGE}
\Css{h1.titleHead { font-weight: normal; }}
\web@headsize{.sectionHead, .likesectionHead}{\Large}
\web@headsize{.subsectionHead, .likesubsectionHead}{\large}
\web@headsize{.subsubsectionHead, .likesubsubsectionHead}{\normalsize}
% a column as wide as the PDF's, in ems of the text, so lines hold about as
% much; narrower screens get the screen's width
\web@Css{html body { max-width: \web@em\textwidth;
    margin-block-start: 2.5em; margin-block-end: 4em; }}
% display math sits in a <table>, whose cell spacing (2px) MathJax's own
% tables inherit; that pulls apart fractions, limits and matrices
\Css{mjx-container { border-spacing: 0; }}
% \intertext, which tables.lua takes out of the alignment
\Css{.intertext { display: block; margin-block: 0.5em; }}
% numbered equations: tex4ht puts the number in a cell of its own, which
% pushes the equation left of center; center it on the line, as TeX does,
% with the number at the right edge
\Css{table.equation, table.equation > tbody { display: block; }}
\Css{table.equation tr { display: grid; grid-template-columns: 1fr minmax(0, auto) 1fr; align-items: center; }}
\Css{table.equation td { grid-column: 2; min-width: 0; padding: 0; }}
\Css{table.equation td.eq-no { grid-column: 3; justify-self: end; }}
% pictures: nothing wider than the column; a picture in display math
% (\[ \begin{tikzcd} .. \]) gets the margin MathJax gives display math
\Css{img { max-width: 100\%; height: auto; }}
\Css{figure.figure { margin-block: 1.75em; }}
\Css{table.equation img, table.equation-star img { display: block; margin: 0.7em auto; }}

\makeatother
